\title{xLSTM: Extended Long Short-Term Memory}

\begin{abstract}
The introduction of the constant error carousel and gating mechanisms during the 1990s established the foundational principles behind Long Short-Term Memory (LSTM) networks. Over subsequent decades, LSTMs have maintained their relevance and fueled many achievements in deep learning—serving as the backbone of some of the earliest Large Language Models (LLMs). Yet, the rise of Transformer models, with their inherently parallelizable self-attention mechanism, ushered in a paradigm shift that eclipsed LSTMs, especially at larger scales. In this context, we revisit the potential of LSTMs, probing the extent to which language modeling can advance when these networks are expanded to billions of parameters, enhanced with contemporary LLM advancements, and their longstanding weaknesses addressed. To this end, we first propose exponential gating, supplemented by tailored normalization and stabilization strategies. Next, we overhaul the memory component of LSTMs, resulting in two new variants: (i) the sLSTM, featuring a scalar memory with scalar updates and a novel approach to memory mixing, and (ii) the mLSTM, which supports full parallelization through matrix memory and a covariance-based update mechanism. By embedding these enhanced LSTM modules within residual block frameworks, we develop xLSTM blocks that can be stacked to construct comprehensive xLSTM architectures. These improvements—incorporating exponential gating and restructured memory—substantially advance the expressivity and performance of xLSTM, enabling them to compete favorably against leading Transformer and State Space Models with respect to both accuracy and scalability.\\
Code available at: \url{https://github.com/NX-AI/xlstm}
\end{abstract}

\section{Introduction}
\label{sec:intro}

Concepts from Long Short-Term Memory (LSTM)~\citep{Hochreiter:91,Hochreiter:97nips,Hochreiter:97}, notably the constant error carousel and the gating mechanisms, were developed to address the issue of vanishing gradients in recurrent neural networks~\citep{Hochreiter:91,Hochreiter:00book}:

\begin{align}
\label{eq:lstm_idea}
  c_t \ = \ f_t \ c_{t-1} \ + \ 
          i_t \  z_t \ , \quad  
          h_t \ = \ o_t \ \psi({c_t}) \ . 
\end{align}

The constant error carousel operates through an additive modification of the cell state $c_{t-1}$ (illustrated in green), incorporating cell inputs $z_t$, with this process being regulated by sigmoid gates (depicted in blue). The modification is governed by the input gate $i_t$ and the forget gate $f_t$, which together modulate the cell state's update, whereas the output of the memory cell—the hidden state $h_t$—is determined by the output gate $o_t$. Before gating, the cell state is subjected to a normalization or squashing function $\psi$, after which output gating generates the hidden state.

LSTMs have achieved notable success across a range of fields \citep{Hochreiter:01,Hochreiter:07,Schmidhuber:15}, and dominated text generation approaches until the introduction of Transformers in 2017~\citep{Vaswani:17}. Their utility has been proven in various sequence processing challenges, including text production~\citep{Graves:13,Karpathy:15blog}, handwriting synthesis~\citep{Graves:13}, sequence-to-sequence translation~\citep{Sutskever:14nips}, program evaluation~\citep{Zaremba:14arxiv}, image captioning~\citep{Karpathy:15,Hossain:19}, source code generation~\citep{Karpathy:15blog}, rainfall-runoff prediction~\citep{Kratzert:18,Kratzert:19}, and developing hydrological models for flood alerts~\citep{Nearing:24}. Within reinforcement learning, LSTMs consistently outperform other sequence modeling architectures, as evidenced by their role in the AlphaStar system for StarCraft II~\citep{Vinyals:17short}, the OpenAI Five framework for Dota 2~\citep{OpenAI:19}, and controllers in nuclear fusion reactors~\citep{Degrave:22short}. LSTMs are particularly adept at learning high-level abstractions, skillfully capturing and retaining semantic patterns within their memory structures~\citep{Karpathy:15blog}, as highlighted by discoveries such as neurons encoding numbers and syntax~\citep{Lakretz:19}, linguistic features~\citep{Bau:19}, and sentiment~\citep{Radford:17arxiv}. Their continued deployment in impactful real-world systems~\citep{Degrave:22short,Nearing:24} attests to their enduring relevance and reliability.

Although LSTMs have achieved notable accomplishments, they exhibit three primary shortcomings: (i) Difficulty in revising memory contents once stored. This issue is illustrated through the {\it Nearest Neighbor Search} problem (additional details are available in Appendix~\ref{sec:appNearestNeighborSearch}): Given a reference vector, the model is required to process a sequence sequentially in order to retrieve the most similar vector’s associated value at the conclusion of the sequence. The left panel of Figure~\ref{fig:lstmProblems} depicts the mean squared error associated with this task. Standard LSTMs are unable to update a previously stored value if a vector with greater similarity is encountered later in the sequence, whereas the proposed xLSTM addresses this via the use of exponential gating. (ii) Constrained memory capacity, as information must be encoded into scalar-valued cell states. This weakness is highlighted using the {\it Rare Token Prediction} task. In the right panel of Figure~\ref{fig:lstmProblems}, perplexity scores on the Wikitext-103~\citep{Merity:17} dataset are plotted according to token frequency. The performance of LSTM degrades for infrequent tokens due to restricted capacity for memory storage. Our xLSTM approach overcomes this through the implementation of a matrix-based memory. (iii) Inherent sequential dependency stemming from recurrent hidden-to-hidden transitions, which inhibits parallel computation and results from the entanglement of memory across time steps.

The shortcomings of LSTM have led to the rise of Transformers~\citep{Vaswani:17} in the field of language modeling. This prompts the question: by addressing these drawbacks and scaling LSTMs to the scale of present-day Large Language Models, what levels of performance in language modeling can be attained?

\section{Extended Long Short-Term Memory}
\label{sec:xLSTM}

In addressing the shortcomings of LSTM networks, Extended Long Short-Term Memory (xLSTM) incorporates two significant changes to the foundational LSTM formulation in Equation~\eqref{eq:lstm_idea}. The enhancements—exponential gating and innovative memory architectures—expand the LSTM family by introducing: (i) sLSTM (see Section~\ref{sec:sLSTM}), characterized by a scalar-valued memory, scalar updates, and the incorporation of memory mixing; and (ii) mLSTM (see Section~\ref{sec:mLSTM}), which employs a matrix-based memory and updates via a covariance (outer product) mechanism that allows for full parallelization. Both variants utilize exponential gating to improve upon the classical LSTM framework. The mLSTM, to facilitate parallel computation, removes memory mixing—specifically, the recurrent hidden-to-hidden connections. Moreover, both sLSTM and mLSTM can be generalized to support multiple memory cells. In this setting, sLSTM continues to mix memory information across the cells, whereas mLSTM treats multi-cell and multi-head configurations equivalently. sLSTM also supports multiple heads without inter-head memory mixing, but retains cell-wise mixing within each head. Introducing multiple heads into the sLSTM architecture in conjunction with exponential gating results in a novel paradigm for memory mixing. For mLSTM, distinctions between heads and cells are functionally redundant.

By incorporating these novel LSTM variants within residual block modules, we obtain xLSTM blocks (refer to Section~\ref{sec:xLSTMblock}). Building architectures by stacking these xLSTM blocks with residual connections leads to what we call xLSTM architectures (see Section~\ref{sec:xLSTMarchitecure}). The structure and components of the xLSTM architecture are depicted in Figure~\ref{fig:xlstm_sketch}.

\subsection{Review of the Long Short-Term Memory}
\label{sec:orgLSTM}

The foundational LSTM model~\citep{Hochreiter:91,Hochreiter:97nips,Hochreiter:97} established the scalar memory cell as its core unit for information processing and retention, effectively mitigating the vanishing gradient issue~\citep{Hochreiter:91,Hochreiter:00book} via the mechanism known as the constant error carousel, or the cell state update. This memory cell is governed by three distinct gates: the input gate, the output gate, and the forget gate, the last of which was incorporated by \citet{Gers:00}. 
The LSTM memory cell's update equations at a given time $t$ are:

\begin{align}
c_t \ &= \  f_t \ c_{t-1} \ + \ i_t \ z_t &  & &\text{cell state} \\
h_t  \ &= \ o_t \ \tilde{h}_t \ , & \tilde{h}_t \ &= \ \psi \left( c_t\right)
  &\text{hidden state} \\
z_t \ &= \ \varphi \left( \tilde{z}_t \right) \ , 
  &\tilde{z}_t \ &=  \ \mathbf{w}^\top_{z} \ \mathbf{x}_t \ + \
  r_{z}  h_{t-1} \ + \  b_{z} \ \
  &\text{cell input} \\
i_t \ &= \ \sigma \left( \tilde{i}_t  \right) \ , 
  &\tilde{i}_t \ &= \ \mathbf{w}^\top_{i} \ \mathbf{x}_t \ + \
  r_{i}  \ h_{t-1} \ + \  b_{i} \ \
  &\text{input gate} \\
f_t \ &= \ \sigma \left(  \tilde{f}_t \right) \ , 
  &\tilde{f}_t \ &= \ \mathbf{w}^\top_{f} \ \mathbf{x}_t  \ + \
  r_{f}  \ h_{t-1} \ + \  b_{f} \ \
  &\text{forget gate} \\
o_t \ &= \ \sigma \left( \tilde{o}_t \right) \ , 
  &\tilde{o}_t  \ &= \ \mathbf{w}^\top_{o} \ \mathbf{x}_t \ + \
  r_{o}  \ h_{t-1} \ + \  b_{o} \ \
  &\text{output gate} 
\end{align}

The vectors $\mathbf{w}_{z}$, $\mathbf{w}_{i}$, $\mathbf{w}_{f}$, and $\mathbf{w}_{o}$ denote the input weights linking the input $\mathbf{x}_t$ to the cell input, input gate, forget gate, and output gate, respectively. The parameters $r_{z}$, $r_{i}$, $r_{f}$, and $r_{o}$ are the recurrent weights connecting the hidden state $h_{t-1}$ to the cell input, input gate, forget gate, and output gate, in that order. The biases associated with these connections are given by $b_{z}$, $b_{i}$, $b_{f}$, and $b_{o}$. Activation functions for the cell input and the hidden state are denoted as $\varphi$ and $\psi$ (most commonly $\tanh$). The function $\psi$ serves to bound the cell state values, which would otherwise be unconstrained. All gating units employ the sigmoid activation function, that is, $\sigma \left( x \right) = 1/(1+\exp(-x))$. Later variations introduced vectors of scalar memory cells, $\mathbf{c}_t \in \mathbb{R}^{d}$, enabling the incorporation of recurrent weight matrices $\mathbf{R} \in \mathbb{R}^{d \times d}$, which facilitates mixing among outputs from distinct memory cells~\citep{Greff:15}; further explanation is provided in Appendix~\ref{sec:apporgLSTM}. Experiments involving removal of these elements demonstrated the essential role of all components in the memory cell architecture~\citep{Greff:15}.

\subsection{sLSTM}
\label{sec:sLSTM}

To enable LSTMs to modify their storage choices, we add exponential gates (shown in red) along with mechanisms for normalization and stabilization. Specifically, we utilize exponential activation functions for both the input and forget gates. For the normalization process, we define a normalizer state, which accumulates the product of the input gate and all subsequent forget gates.

The forward pass for the scalar sLSTM is:

\begin{align}
\label{eq:slstmforward}
c_t \ &= \  f_t \ c_{t-1} \ + \ i_t \ z_t & & 
 &\text{cell state} \\
n_t \ &= \  f_t \ n_{t-1} \ + \ i_t  & & 
 &\text{normalizer state} \\
h_t \ &= \ o_t \ \tilde{h}_t \ ,
  &\tilde{h}_t \ &= \ c_t / n_t
&\text{hidden state} \\
z_t \ &= \ \varphi \left( \tilde{z}_t \right) \ , 
  &\tilde{z}_t \ &=  \ \mathbf{w}^\top_{z} \ \mathbf{x}_t \ + \
  r_{z}  h_{t-1} \ + \  b_{z}
  &\text{cell input} \\
i_t \ &= \ \!\exp \left( \tilde{i}_t  \right) \ , 
  &\tilde{i}_t \ &= \ \mathbf{w}^\top_{i} \ \mathbf{x}_t \ + \
  r_{i}  \ h_{t-1} \ + \  b_{i}
  &\text{input gate} \\
f_t \ &= \ \sigma \left(  \tilde{f}_t \right) \ \text{OR} \
\!\exp \left(  \tilde{f}_t \right) \ ,
  &\tilde{f}_t \ &= \ \mathbf{w}^\top_{f} \ \mathbf{x}_t  \ + \
  r_{f}  \ h_{t-1} \ + \  b_{f}
  &\text{forget gate} \\
o_t \ &= \ \sigma \left( \tilde{o}_t \right) \ , 
  &\tilde{o}_t  \ &= \ \mathbf{w}^\top_{o} \ \mathbf{x}_t \ + \
  r_{o}  \ h_{t-1} \ + \  b_{o}
  &\text{output gate} 
\end{align}

We adapt the initial LSTM gating strategies—specifically, gating mechanisms that depend on the input and/or hidden state, along with a bias term—to the new model designs. Since exponential activation functions may generate large outputs that risk numerical overflow, we enhance gate stability by introducing an auxiliary state, \mbox{$m_t$~\citep{Milakov:18arxiv}}:

\begin{align}
\label{eq:slstmstabil}
m_t \ &= \ \max \left( \log ( f_t ) + m_{t-1} , \log ( i_t ) \right) &\text{stabilizer state} \\
i'_t \ &= \ \!\exp \left( \log \left ( i_t \right) - m_t \right) \ = \ \!\exp \left( \tilde{i}_t - m_t \right) \ \, 
  &\text{stabil. input gate} \\
f'_t \ &= \ \!\exp \left( \log \left( f_t \right) + m_{t-1} - m_t \right) \ \, 
  &\text{stabil. forget gate}
\end{align}

As demonstrated in Appendix~\ref{sec:appsLSTM}, substituting $f_t$ with $f'_t$ and $i_t$ with $i'_t$ during the forward pass leaves both the overall network output and the gradients of the loss relative to the parameters unaffected.

\paragraph{New Memory Mixing.} Similar to the standard LSTM (refer to Appendix~\ref{sec:appsLSTM}), sLSTM allows for the use of several memory cells. The presence of multiple memory cells facilitates mixing through the recurrent connections $\mathbf{R}_{\mathbf{z}}$, $\mathbf{R}_{\mathbf{i}}$, $\mathbf{R}_{\mathbf{f}}$, and $\mathbf{R}_{\mathbf{o}}$, linking the hidden state $\mathbf{h}$ to the memory cell input $\mathbf{z}$ as well as to the gates $\mathbf{i}$, $\mathbf{f}$, and $\mathbf{o}$. What distinguishes the current approach is the inclusion of exponential gating within the process of memory mixing. Furthermore, sLSTM can support several heads, permitting internal memory mixing among the cells within a single head but not between different heads. The integration of head structures and exponential gating in sLSTM introduces a novel mechanism for memory mixing.

\subsection{mLSTM}
\label{sec:mLSTM}

In order to expand the memory capacity of LSTMs, we generalize the traditional scalar cell state $c \in \mathbb{R}$ to a matrix form $\mathbf{C} \in \mathbb{R}^{d \times d}$. This modification allows retrieval operations to be executed using matrix multiplication. At each time step $t$, a pair of vectors—the key $\mathbf{k}_t \in \mathbb{R}^d$ and the value $\mathbf{v}_t \in \mathbb{R}^d$ (as per Transformer nomenclature)—are stored in memory. Subsequently, at a future time $t + \tau$, a query vector $\mathbf{q}_{t+\tau} \in \mathbb{R}^d$ is employed to recall the previously stored value $\mathbf{v}_t$. This approach corresponds to the framework established by Bidirectional Associative Memories (BAMs) \citep{Kohonen:72,Anderson:72,Nakano:72,Anderson:77}.

To encode a key-value pair, the covariance update rule~\citep{Sejnowski:77,Dayan:91} is utilized.

\begin{align}
\mathbf{C}_t \ &= \ \mathbf{C}_{t-1} \ + \ \mathbf{v}_{t} \ \mathbf{k}_t^\top \ .
\end{align}

We posit that a layer normalization is performed prior to mapping the inputs to keys and values, resulting in these vectors possessing zero mean. The covariance update mechanism achieves optimality~\citep{Dayan:91} with respect to maximizing the separability of retrieved binary codes, which corresponds to optimizing the signal-to-noise ratio. Enhanced separability may be achieved by constraining retrieval to pairwise interactions, thereby accepting quadratic computational complexity, as is the case with attention mechanisms~\citep{Krotov:16,Krotov:17,Ramsauer:21}. Notably, the covariance update procedure aligns with the formulation of Fast Weight Programmers~\citep{Schmidhuber:92ncfastweights,Schlag:21}, which have subsequently incorporated a fixed decay factor applied to $\mathbf{C}_{t-1}$ as well as a constant learning rate for 
$\mathbf{v}_{t} \mathbf{k}_t ^\top$~\citep{Ba:16}.
Motivated by this, we adapt the covariance update methodology within the LSTM architecture, mapping the forget gate to the decay parameter, the input gate to the learning rate, and utilizing the output gate for modulating the recovered vector.

Within this matrix memory framework, the normalizer state is computed as a weighted aggregation of key vectors, where each key is scaled by the product of its associated input gate and the sequence of subsequent forget gates. This normalizer state serves to track the cumulative effect of the gating mechanisms. Given that the dot product between the query and the normalizer state may approach zero, we follow prior work~\citep{Sun:23arxiv} in taking its absolute value and ensuring it does not fall below a specified threshold (commonly set to 1.0). The forward computation for mLSTM is as follows:

\begin{align}
\label{eq:mlstm_recurrent_begin}
\mathbf{C}_t \ &= \  f_t \ \mathbf{C}_{t-1} \ + \ 
  i_t \ \mathbf{v}_t \ \mathbf{k}_t^\top & &  &\text{cell state} \\
\mathbf{n}_t \ &= \  f_t \ \mathbf{n}_{t-1} \ + \ 
  i_t \ \mathbf{k}_t & &  &\text{normalizer state} \\
\mathbf{h}_t  \ &= \ \mathbf{o}_t \ \odot \ \tilde{\mathbf{h}}_t
\ , \qquad \qquad 
\tilde{\mathbf{h}}_t \ = \   \mathbf{C}_t \mathbf{q}_t \ / \ 
  \max \left\{ |\mathbf{n}_t^\top \mathbf{q}_t|, 1 \right\} 
   &&&\text{hidden state} \\
\mathbf{q}_t \ &= \ \mathbf{W}_q \ \mathbf{x}_t \ + \ \mathbf{b}_q  & &  &\text{query input} \\
\mathbf{k}_t \ &= \ \frac{1}{\sqrt{d}} \mathbf{W}_k \ \mathbf{x}_t \ + \ \mathbf{b}_k  & &  &\text{key input} \\
\mathbf{v}_t \ &= \ \mathbf{W}_v \ \mathbf{x}_t \ + \ \mathbf{b}_v  & &  &\text{value input} \\
i_t \ &= \ \!\exp \!  \! \left( \tilde{i}_t  \right) \ , 
 \qquad \qquad \ \ \ \ \,
  \tilde{i}_t \ = \ \mathbf{w}^\top_{i} \ \mathbf{x}_t \ + \  b_{i} \ \ \ 
  &&&\text{input gate} \\
f_t \ &= \ \sigma \!  \left(  \tilde{f}_t \right) \ \text{OR} \ \!\exp \!  \! \left(  \tilde{f}_t \right) , \ \ \: \!
\tilde{f}_t \ = \ \mathbf{w}^\top_{f} \ \mathbf{x}_t  \ + \
  b_{f}
  &&& \text{forget gate} \\
\label{eq:mlstm_recurrent_end}
\mathbf{o}_t \ &= \ \sigma \left( \tilde{\mathbf{o}}_t \right) \ , \qquad \qquad \quad \ \ \ \,
  \tilde{\mathbf{o}}_t  \ = \ \mathbf{W}_{\mathbf{o}} \ \mathbf{x}_t \ + \
  \mathbf{b}_{\mathbf{o}}
  &&&\text{output gate} 
\end{align}

mLSTM, similar to the standard LSTM, may utilize several memory cells. In the context of mLSTM, using multiple heads or multiple cells leads to the same result, since no memory mixing takes place. To regulate the exponential gating mechanisms within mLSTM, we adopt the stabilization strategies applied in sLSTM (refer to Equation~\ref{eq:slstmstabil}). Because there is no mixing of memory in mLSTM, its recurrence relation can be restructured into a parallelizable form. Additional information is provided in Appendix~\ref{sec:appmLSTM}.

\subsection{xLSTM Architecture}
\label{sec:xLSTMarch}

\paragraph{xLSTM Blocks.}
\label{sec:xLSTMblock}
An xLSTM block is designed to generate a non-linear representation of preceding information within a high-dimensional feature space, thereby enhancing the distinction between various historical contexts. Effectively distinguishing different pasts is crucial for accurate next-step predictions, such as forecasting the upcoming token in a sequence. This design is motivated by Cover’s Theorem~\citep{Cover:65}, which posits that non-linearly mapped patterns are more likely to be linearly separable in a higher-dimensional space compared to the original. We investigate two types of residual block structures: (i) A residual block with up-projection applied after non-linear summarization (analogous to Transformers), where the history is first non-linearly aggregated in the original feature space, subsequently projected linearly into a higher-dimensional space, transformed by a non-linear activation, and finally mapped linearly back to the original dimension; refer to the left panel of Figure~\ref{fig:Backbones} and the third column in Figure~\ref{fig:xlstm_sketch}. A comprehensive schematic is shown in Figure~\ref{fig:appsLSTM_detailed} in the appendix. (ii) A residual block with an up-projection preceding non-linear operations (inspired by State Space Models), which first linearly projects to a higher-dimensional representation,

It performs a nonlinear aggregation of previous information within a high-dimensional space before applying a linear transformation to return to the original representation. In the case of an xLSTM block that incorporates an sLSTM, the post up-projection block is utilized. Conversely, for an xLSTM block equipped with an mLSTM, the pre up-projection block is preferred, owing to the increased memory capacity available in the high-dimensional space. For a visual illustration, refer to the left panel of Figure~\ref{fig:Backbones}, the third column of Figure~\ref{fig:xlstm_sketch}, or consult Figure~\ref{fig:appsLSTM_detailed} in the appendix for further clarification.

\paragraph{xLSTM Architecture. }
\label{sec:xLSTMarchitecure}
The xLSTM model is developed by applying residual stacking of fundamental modules~\citep{Srivastava:15,He:16}. For the backbone, we adopt the prevalent pre-LayerNorm~\citep{Ba:16b} residual structure, which is standard in current Large Language Models. This configuration is illustrated in the final column of Figure~\ref{fig:xlstm_sketch}.

\subsection{Memory and Speed Considerations}

In contrast to Transformers, xLSTM networks exhibit linear computational complexity and maintain constant memory usage as the sequence length increases. Due to the compressive nature of xLSTM memory, these networks are particularly advantageous for use in industrial settings and for deployment on edge devices.

While mLSTM’s memory mechanism does not depend on additional parameters, its reliance on $d \times d$ matrix storage and $d \times d$ matrix operations incurs significant computational overhead. Our approach deliberately balances the increased memory capacity with this computational burden. However, since these matrix operations are well-suited for parallel execution on GPUs, the actual increase in wall clock time remains minimal.

Although mLSTM supports parallel processing similar to FlashAttention~\citep{Dao:22,Dao:23} and GLA~\citep{Yang:23arxiv}, sLSTM lacks such parallelizability because of its memory mixing mechanism (hidden-hidden dependencies). Nonetheless, we introduced an efficient CUDA-based implementation that leverages GPU memory optimizations down to the register level, resulting in performance that is usually less than twice as slow as mLSTM.

\section{Related Work}
\label{sec:related}

\paragraph{Linear Attention.} A variety of approaches have been proposed to address the quadratic scaling of Transformers with respect to context length, aiming to achieve attention with linear time complexity. The Synthesizer sidesteps direct token--token interactions by learning synthetic attention scores~\citep{Tay:20arxiv}. Linformer achieves a linearized form of self-attention via low-rank projections of the attention matrix~\citep{Wang:20arxiv}. The Linear Transformer implements a linear variant of the attention mechanism \citep{Katharopoulos:20}. The Performer uses a method based on positive orthogonal random features to produce a linear approximation of the softmax attention~\citep{Choromanski:21}. Alternatives to attention, such as fast, long-sequence convolutions, have been introduced in models like the Structured Global Convolution (SGConv)~\citep{Li:22} and the Hyena Hierarchy~\citep{Poli:23}.

\paragraph{State Space Models.} In recent times, State Space Models (SSMs) have garnered significant attention due to their linear computational complexity with respect to context length and competitive results in comparison with Transformers. Early notable models in this area include the Structured State Space sequence model (S4)~\citep{Gu:21}. Subsequent advancements have led to the introduction of the Diagonal State Space (DSS) model~\citep{Gupta:22}, Gated State Space (GSS) models~\citep{Mehta:22}, S5 model~\citep{Smith:22}, Bidirectional Gated SSM (BiGS)~\citep{Wang:22}, H3 model~\citep{Fu:23}, and Mamba~\citep{Gu:24arxiv}.

\paragraph{Recurrent Neural Networks.} Recurrent Neural Networks (RNNs) have been proposed as alternatives to Transformers and attention mechanisms, primarily because their computational cost scales linearly with sequence length. Recent advancements, such as RNNs equipped with Deep Linear Recurrent Units (LRUs), have achieved notable success in language modeling tasks~\citep{Orvieto:23, De:24arxiv}. Additionally, models like the Hierarchically Gated Linear RNN (HGRN)~\citep{Qin:23} and its successor, HGRN2~\citep{Qin:24arxiv}, have demonstrated strong performance. Among RNN-based methods for large-scale language modeling, RWKV~\citep{Peng:23arxivshort,Peng:24arxivshort} is recognized for delivering results comparable to Transformer models.

\paragraph{Gating.}
The concept of gating, originally central to LSTM architectures, has subsequently been revisited and reformulated in numerous modern methods. For instance, gating mechanisms are incorporated in HGRN~\citep{Qin:23}, HGRN2~\citep{Qin:24arxiv}, Gated Linear Attention (GLA)~\citep{Yang:23arxiv}, Gated State Space (GSS) models~\citep{Mehta:22}, Bidirectional Gated SSM (BiGS)~\citep{Wang:22}, Moving Average Equipped Gated Attention (MEGA)~\citep{Ma:22}, RWKV~\citep{Peng:23arxivshort}, and Mamba~\citep{Gu:24arxiv}.

\paragraph{Covariance Update Rule.} In order to improve memory capacity, we integrated a matrix memory governed by a covariance update rule into the mLSTM cell. Related approaches employing this update paradigm include Fast Weight Programmers \citep{Schmidhuber:92ncfastweights,Schlag:21}, RWKV-5 and RWKV-6~\citep{Peng:24arxivshort}, Retention~\citep{Sun:23arxiv}, Linear Transformer~\citep{Katharopoulos:20}, and HGRN2~\citep{Qin:24arxiv}.

\paragraph{Most Related.} The models that are most akin to xLSTM in their conceptual design include Retention~\citep{Sun:23arxiv}, RWKV~\citep{Peng:23arxivshort,Peng:24arxivshort}, and HGRN2~\citep{Qin:24arxiv}. These architectures employ mechanisms such as matrix-based memory and gating. Nonetheless, unlike the recently introduced sLSTM, these existing methods do not support memory mixing. The capability for memory mixing is critical for addressing state tracking tasks, making LSTMs inherently more expressive compared to State Space Models (SSMs) and Transformers~\citep{Merrill:24,Deletang:23}. Effective state tracking is essential for tasks such as code evaluation or maintaining the coherence of entities across lengthy texts.

\paragraph{Residually Stacking Architectures.} Consistent with most state-of-the-art deep learning frameworks, xLSTM models utilize residual stacking of modular components~\citep{Srivastava:15,He:16}.

This architectural strategy paved the way for the advancements in deep convolutional neural networks~\citep{He:16} and also forms the backbone of Transformer-based systems~\citep{Vaswani:17}. Transformers, in turn, underpin modern Large Language Models (LLMs) such as GPT-3~\citep{Brown:20short}, ChatGPT~\citep{Schulman:22short}, GPT-4~\citep{Achiam:23gpt4short}, Megatron-LM~\citep{Shoeybi:19}, Gopher~\citep{Rae:21short}, ERNIE 3.0 Titan~\citep{Wang:21arxivshort}, GLaM~\citep{Du:21glamshort}, Chinese M6~\citep{Lin:21}, multilingual AlexaTM 20B~\citep{Soltan:22}, OPT~\citep{Zhang:22opt}, Chinchilla~\citep{Hoffmann:22short}, BLOOM~\citep{Scao:22arxivshort}, GLM-130B~\citep{Zeng:22arxivshort}, LaMDA~\citep{Thoppilan:22short}, PaLM~\citep{Chowdhery:22short}, Llama~\citep{Touvron:23arxiv}, and Gemini~\citep{GeminiTeam:23arxiv,Reid:24arxiv}.

\section{Experiments}
\label{sec:Exp}

This section presents an empirical assessment of xLSTM, benchmarking its performance in language modeling against established techniques. Section~\ref{sec:ExpSyntethic} explores the unique abilities of xLSTM through a set of synthetic benchmarks. Section~\ref{sec:ExpComparison} provides a comparative analysis of the validation perplexity achieved by contemporary language models, all trained on 15B SlimPajama tokens~\citep{Soboleva:23}, including a detailed ablation study focused on xLSTM. Additionally, we analyze the scaling patterns of these methods following the frameworks set out by \citet{Kaplan:20} and \citet{Brown:20short}. 

In Section~\ref{sec:ExpLanguage}, we present a more comprehensive study on language modeling, comparing xLSTM and the top-performing baselines identified in Section~\ref{sec:ExpComparison}, now trained on an expanded 300B SlimPajama tokens~\citep{Soboleva:23}. This investigation comprises several components: evaluating the ability of each approach to generalize to longer sequences; measuring validation set perplexity and downstream performance~\citep{Sutawika:24}; testing across 571 domains from the PALOMA language benchmark~\citep{Magnusson:23arxivshort}; and re-examining scaling properties, this time under a twentyfold increase in training data volume.

\label{sec:xlstm_blocks_notation}
Throughout our experiments, we denote xLSTM[$a$:$b$] as representing a proportion of $a$ mLSTM-based xLSTM blocks to $b$ sLSTM-based blocks. For instance, the configuration xLSTM[7:1] signifies that, among eight total blocks, seven utilize mLSTM while one utilizes sLSTM. If the total number of blocks is set to 48, this configuration yields 42 mLSTM-based blocks and 6 sLSTM-based blocks. Additionally, in every experimental setup, we employ up-projection blocks before and after mLSTM and sLSTM blocks, respectively.

\subsection{Synthetic Tasks and Long Range Arena}
\label{sec:ExpSyntethic}
To begin, we evaluate the impact of xLSTM's exponential gating combined with memory mixing on tasks involving formal languages~\citep{Deletang:23}. Next, the capabilities of the novel matrix memory in xLSTM are examined using the Multi-Query Associative Recall task~\citep{Arora:23arxiv}. Lastly, we investigate how well xLSTM handles lengthy input sequences through experiments on the Long Range Arena benchmark~\citep{Tay:21}.

\paragraph{Assessment of Exponential Gating with Memory Mixing in xLSTM.} We evaluate the effectiveness of xLSTM’s exponential gating coupled with memory mixing, a mechanism posited to address state tracking challenges~\citep{Merrill:24, Merrill:23}. Building upon and broadening the formal language benchmarks of \citet{Deletang:23}, we incorporate multi-length training to facilitate the study of length generalization. Further details concerning the experimental setup and expanded findings are presented in Appendix~\ref{sec:appExpSynthetic-formalLang}. We conduct a comparative analysis between xLSTM and a spectrum of architectures such as Transformers, State Space Models, and conventional Recurrent Neural Networks. For each model, task-relevant token accuracy is reported, scaled within the range [0,1], where 0 represents random guessing and 1 denotes flawless performance. In these experiments, 2-block variants are considered, including: xLSTM[0:1] (exclusively sLSTM), xLSTM[1:0] (exclusively mLSTM), xLSTM[1:1], Llama, Mamba, RWKV, Retention, Hyena, LSTM, and LSTM embedded within Transformer blocks (LSTM (Block)). Experimental outcomes are summarized in Figure~\ref{fig:formal-main}. Notably, models such as Transformers or State Space Models lacking memory mixing (and thus the ability for explicit state tracking) fail on tasks involving, for example, regular grammars like the parity benchmark. This observation corroborates prior results indicating the inherent limitations of Transformers and State Space architectures in comparison to RNNs for such tasks~\citep{Merrill:24, Merrill:23, Deletang:23}.

\paragraph{Test of xLSTM's Memory Capacities on Associative Recall Tasks.} In this study, we evaluate the enhanced matrix memory mechanism in xLSTM by examining its ability to store information within the Multi-Query Associative Recall benchmark~\citep{Arora:23arxiv}. Here, each sequence consists of key-value pairs, which are randomly selected from an extensive vocabulary and must be retained for subsequent retrieval. To raise the level of challenge relative to the standard version of the task, we expand the number of key-value pairs up to 256 and extend the sequence context window to as many as 2048 tokens. This setup enables a more comprehensive investigation into the memory capacity of several models. We perform a comparative assessment involving 2-block variants of Llama, Mamba, RWKV-5, RWKV-6, xLSTM[1:1], and xLSTM[1:0], and gauge performance according to the accuracy of pair retrieval. Since Transformer architectures such as Llama are characterized by memory capacity that scales exponentially with respect to coding dimensionality~\citep{Ramsauer:21}, they set the upper benchmark on this task. Performance outcomes appear in Figure~\ref{fig:mqar-main}. Among the non-Transformer competitors, xLSTM[1:1] achieves the highest results, including for model configurations with fewer parameters. Notably, the inclusion of the sLSTM block does not impede, but instead appears to enhance memory retention, an effect most pronounced in the setting with 256 key-value pairs. Further analyses, including extrapolation experiments demonstrating the generalization of xLSTM memory to unseen context lengths, can be found in Appendix~\ref{sec:appExpSynthetic-mqar}.

\paragraph{Evaluation of xLSTM's Ability to Handle Extended Contexts Using Long Range Arena.} In order to evaluate xLSTM's effectiveness with lengthy sequences and extensive contexts, we benchmark various approaches on the Long Range Arena~\citep{Tay:21}. Across all evaluated tasks, xLSTM achieves robust and reliable results, indicating that the xLSTM model is highly adept at managing diverse challenges associated with long context scenarios.

Additional information can be found in Appendix~\ref{sec:appExpSynthetic-lra}.

\subsection{Method Comparison and Ablation Study}
\label{sec:ExpComparison}

In order to investigate our central research question—namely, the capabilities of our novel LSTM variants when scaled for language modeling—we train xLSTMs, Transformers, State Space Models, and several other approaches using 15B tokens from SlimPajama within an identical auto-regressive framework. The performance of all trained models is evaluated on a shared validation set, and we additionally conduct ablation analyses on the xLSTM architectures.

\paragraph{Comparing xLSTM to Other Methods.} To facilitate a fair comparison, all models are trained on a dataset of 15B tokens sourced from SlimPajama~\citep{Soboleva:23}. Their effectiveness is gauged via perplexity computed on the validation data. We benchmark the following architectures: our proposed xLSTM; GPT-3 (Transformer)~\citep{Brown:20short}; Llama (Transformer)~\citep{Touvron:23arxiv}; H3 (SSM)~\citep{Fu:23}; Mamba (SSM)~\citep{Gu:24arxiv}; RWKV-4, RWKV-5, and RWKV-6 (all RNN variants)~\citep{Peng:23arxivshort, Peng:24arxivshort}; GLA (linear Transformer)~\citep{Yang:23arxiv}; HGRN and HGRN2 (RNNs)~\citep{Qin:23, Qin:24arxiv}; RetNet (linear Transformer)~\citep{Sun:23arxiv}; Hyena (linear Transformer)~\citep{Poli:23}; xLSTM[1:0]; and xLSTM[7:1] (as described in Section~\ref{sec:xlstm_blocks_notation}). 

All models utilize mixed precision for training, although, in the case of RWKV-5, RWKV-6, GLA, and HGRN2, the mixed-precision regime does not employ the PyTorch automated mixed precision feature (see Appendix Section~\ref{sec:appGenTrainProc} for further clarification). 

The comparison is organized into (a) Transformers, (b) State Space Models (SSMs), and (c) Recurrent Neural Networks (RNNs), along with linear Transformer variants that replace standard attention with linear operations. All models are parameter-matched to a 350M parameter GPT-3 (i.e., embedding dimension of 1024 and 24 residual layers). Notably, only GPT-3 employs weight sharing for input and output embeddings, which marginally reduces its parameter count. As demonstrated in Table~\ref{tab:spaj15b_model_results}, xLSTM achieves superior validation perplexity compared to all reference models. Additional information can be found in Appendix~\ref{sec:appExpComparison}. The scaling trends shown in Figure~\ref{fig:temp_sclaw15B} suggest that xLSTM continues to deliver strong performance as model size increases.

\paragraph{Ablation Studies.}
As illustrated in Table~\ref{tab:spaj15b_model_results} and Figure~\ref{fig:temp_sclaw15B}, xLSTM demonstrates strong language modeling performance when trained on 15B tokens from the SlimPajama dataset. To systematically analyze the impact of architectural modifications from LSTM to xLSTM, we progressively adapt a standard LSTM framework towards xLSTM in several stages. The initial step involves embedding LSTM layers within pre-LayerNorm residual architectures. Subsequently, we transition to utilizing a post up-projection block. The final modification incorporates both exponential gating and the matrix memory mechanism.

Table~\ref{tab:ablstudies} (top) presents the findings.

The ablation experiments indicate that both exponential gating and matrix memory are key contributors to the observed performance gains. Given the central role of gating in RNNs and State Space Models, we also examine alternative gating approaches. As summarized in Table~\ref{tab:ablstudies} (bottom), making each gate trainable and responsive to input leads to further incremental improvements. Further analyses concerning the backbone components are provided in Appendix~\ref{sec:appAblDetails}.

\subsection{xLSTM as Large Language Model}
\label{sec:ExpLanguage}

This investigation concludes with extensive language modeling experiments, wherein we assess xLSTM’s capability as a large language model (LLM). To this end, we scale up the training dataset, utilizing 300B tokens sourced from SlimPajama. Notably, this volume of training data is consistent with that employed by models such as Mamba~\citep{Gu:24arxiv} and Griffin~\citep{De:24arxiv}.

We benchmark xLSTM against RWKV-4, Llama, and Mamba, each chosen to represent a distinct method class described in Section~\ref{sec:ExpComparison}. RWKV-4 is included as the RNN baseline since stable training precision settings for RWKV-5, RWKV-6, and HGRN2 became available only after the commencement of the 300B token training experiments (refer to Appendix~\ref{sec:appGenTrainProc} for details). Our study covers a range of model sizes (125M, 350M, 760M, 1.3B), examining each model’s ability to extrapolate to longer sequence lengths and analyzing validation set outcomes. Additionally, we evaluate the models' performance on downstream tasks, their language modeling proficiency across the 571 textual domains provided by the PALOMA benchmark, and explore how their results align with established scaling laws.

\paragraph{Sequence Length Extrapolation.}
\label{sec:spaj300B_ctx_extrapolation}
We begin by evaluating how well 1.3B parameter models—including xLSTM, RWKV-4, Llama, and Mamba—generalize to longer sequence lengths. The training was conducted with a context window of 2048 tokens, but during testing, models are assessed with context lengths extended up to 16384. Results are displayed in Figure~\ref{fig:spaj300B_seq_extrapolation}. Unlike the other approaches, xLSTM consistently preserves low perplexity when exposed to significantly increased context lengths.

\paragraph{Validation Perplexity and Downstream Tasks.}
\label{sec:spaj_downstream_eval}
Next, across every model scale, we assess xLSTM, RWKV-4, Llama, and Mamba by their perplexity on the SlimPajama validation set for next-token prediction, as well as by their effectiveness on downstream benchmarks that evaluate common sense reasoning capabilities.

The third column in Table~\ref{tab:lmeval} presents the validation set perplexity for the various methods considered. For all evaluated model sizes, both xLSTM[1:0] and xLSTM[7:1] achieve the lowest perplexity on the validation set, thus outperforming other approaches in this regard. Performance results on downstream tasks are shown in the remaining columns of Table~\ref{tab:lmeval}. Across nearly every downstream task and model size, xLSTM consistently delivers superior results, with the exception of the ARC task, where Mamba sometimes obtains the highest scores. Additional information can be found in Appendix~\ref{sec:appExpLanguage}.

\paragraph{Performance on PALOMA Language Tasks.} Thirdly, we evaluate models of varying scales—xLSTM, RWKV-4, Llama, and Mamba—on their ability to predict the next token in the PALOMA language tasks~\citep{Magnusson:23arxivshort}. This evaluation is conducted using perplexity as a metric over 571 distinct text domains, spanning sources such as nytimes.com and Reddit's r/depression. Table~\ref{tab:ppleval} presents the results, distinguishing between language modeling tasks (first seven columns) and detailed domain-specific benchmarks (last five columns). On these tasks, xLSTM[1:0] demonstrates superior performance compared to xLSTM[7:1].

xLSTM[1:0] achieves lower perplexity compared to Mamba on 568 of 571 text domains (99.5\%), surpasses Llama in 486 out of 571 cases (85.1\%), and outperforms RWKV-4 in 570 of 571 domains (99.8\%). Further details are provided in Appendix~\ref{sec:appExpLanguage}.

\paragraph{Scaling Laws.} Fourth, we investigate the scaling behavior described by power laws, which can be used to predict the performance of models as their size increases~\citep{Kaplan:20,Brown:20short}. Figure~\ref{fig:temp_sclaw300B} illustrates this scaling pattern. Although all evaluated models exhibit comparable scaling trends, they differ in their respective offsets. Among the models, RWKV-4 demonstrates the poorest performance, with Llama and Mamba following. Notably, xLSTM consistently outperforms Mamba, with the performance gap between xLSTM and Mamba being similar to that observed between Mamba and Llama. The observed scaling characteristics suggest that xLSTM is likely to maintain, or even strengthen, its performance advantage over both Transformer and State-Space architectures as model size grows.

\textbf{Generation Times and Maximal Throughput.} We assess the speed of text generation in Figure~\ref{fig:inference_time_generation} and analyze the peak decoding throughput in Figure~\ref{fig:inference_time_decoding_throughput} (left) across the various xLSTM configurations, each at a 1.3B parameter count. Our evaluation includes baseline comparisons against comparably sized models—Mamba, Llama, and RWKV—as provided by HuggingFace, noting that for Llama, a static key-value cache is enabled. At the moment these tests were conducted, compiling the complete key-value cache for the Transformer, as well as compiling the Mamba model using \texttt{torch.compile}, was not feasible. For generation speed measurements, all models utilize a batch size of 1 with a pre-fill of 16 tokens, a scenario that is particularly advantageous for the Transformer architecture.

The results in Figure~\ref{fig:inference_time_generation} highlight that xLSTM, along with the other recurrent models (Mamba and RWKV-4), exhibit linear complexity with respect to sequence length, whereas Llama demonstrates quadratic scaling. In the throughput experiments, we explore various batch sizes and prefill lengths for Llama. Figure~\ref{fig:inference_time_decoding_throughput} (right) demonstrates that xLSTM accommodates substantially larger batch sizes relative to Llama, attributed to its fixed memory requirements, resulting in superior overall throughput.

\section{Limitations}

(i) Unlike mLSTM, the sLSTM’s approach to memory mixing introduces dependencies that prevent operations from being easily parallelized, which impedes the development of highly parallel, fast implementations. Despite this limitation, we engineered an efficient CUDA kernel for sLSTM, achieving performance at less than twice the runtime of our optimized parallel mLSTM implementation. (ii) It should be noted that the existing CUDA kernels for mLSTM lack significant optimization, resulting in runtimes roughly four times slower than FlashAttention or Mamba’s scan mechanism. There is potential for performance improvements through advanced kernel design similar to FlashAttention. (iii) The matrix memory in mLSTM imposes substantial computational demand, as it necessitates manipulation of $d \times d$ matrices. Nevertheless, because both memory update and retrieval are parameter-free and benefit from matrix operation parallelism, the practical increase in wall-clock time remains modest. (iv) Careful attention is required when setting the initial values for the forget gates. (v) The matrix memory is decoupled from sequence length, but if the sequence length increases substantially, it could exceed the memory’s representational capacity for large contexts. Nonetheless, this has not posed an issue for context sizes up to 16k tokens; see Section~\ref{sec:spaj300B_ctx_extrapolation}. (vi) Owing to the significant computational costs associated with large-scale language model experiments, we did not exhaustively tune the architecture or hyperparameters, especially with larger xLSTM models. We expect that comprehensive optimization will be necessary for xLSTM models to fully realize their capabilities.

\section{Conclusion}
We have provided a partial response to our guiding inquiry: To what extent does language modeling benefit from scaling LSTM architectures to the order of billions of parameters? Based on our findings, the answer is: ``At minimum, scaling LSTM yields results comparable to prevailing models such as Transformers and State Space Models''. Our contributions include augmenting LSTM into xLSTM through the introduction of exponential gating with memory mixing, as well as implementing an alternative memory framework. Empirical results show that xLSTM achieves competitive or superior language modeling performance relative to leading techniques including Transformers and State Space Models. Analysis of the scaling behavior suggests that with increased model size, xLSTM stands as a formidable challenger to Large Language Models reliant on Transformer architectures. Furthermore, xLSTM holds promise for significant influence in related domains such as Reinforcement Learning, Time Series Prediction, and modeling of physical phenomena.

\end{document}